\section{Why the discrete and continuous constants differ}
\label{a:sec:compare}

The two constants are
\begin{equation*}
\Ac(p) = \frac{1-2p}{1-p}\quad\text{(continuous time)},
\qquad
A(p) = \lim_{n\to\infty}\frac{S_n}{(2p)^n}\quad\text{(discrete time)},
\end{equation*}
and the observation recorded in \ChM{} was that the first starts at $1$ where the
second starts at $\tfrac12$, yet halving the first does not superpose the curves
(\cref{a:fig:dtct}).

\begin{figure}[ht]
\centering
\includegraphics[width=0.5\textwidth]{dtctA.png}
\caption{The continuous-time constant $\Ac(p) = (1-2p)/(1-p)$ (green) against
high-precision numerical values of the discrete-time constant $A(p)$ (red points).
Both vanish at $p=\tfrac12$, but at $p=0$ the continuous curve starts at $1$ and
the discrete one at $\tfrac12$, and the two are not related by a constant factor in
between. The black line is a formula obtained by symbolic regression
(\cref{a:sec:GA}); it tracks the points closely but is not exact.}
\label{a:fig:dtct}
\end{figure}

What holds is not $A=\tfrac12\Ac$ but the strict inequality $A>\tfrac12\Ac$ of
\cref{a:eq:halfAc}, and the inequality is attained only in the limit $p\to0$.
Combining \cref{a:eq:Aasym,a:eq:Acts}, the ratio runs monotonically across the
interval:
\begin{equation}
\label{a:eq:ratiolimits}
\frac{A(p)}{\Ac(p)} \;\longrightarrow\; \tfrac12 \quad(p\to0),
\qquad
\frac{A(p)}{\Ac(p)} \;\longrightarrow\; 1 \quad(p\to\tfrac12^-),
\end{equation}
taking the values $0.503$, $0.528$, $0.619$, $0.798$, $0.928$ and $0.988$ at
$p=0.01,\,0.1,\,0.3,\,0.45,\,0.49,\,0.499$.

At $p\to0$ the factor of two is a counting effect. In the binary discrete model an
individual leaves either nothing or two offspring, so a surviving population is
always even (\cref{a:prop:parity}); deep in the subcritical regime, conditioning on
survival means conditioning on a population of exactly two, and the
quasi-stationary mean is $2$. In the continuous-time model a survivor is a single
individual that has not yet died, so the quasi-stationary mean is $1$. The
reciprocals are $\tfrac12$ and $1$.

At $p\to\tfrac12^-$ the factor disappears, because both constants are governed by
the same critical scaling. Expanding \cref{a:eq:Acts} near $p=\tfrac12$ gives
$\Ac(p) = 2\varepsilon/(1+\varepsilon)\to2\varepsilon$, which is precisely
\cref{a:eq:Aasym}. Near criticality the discreteness of the generational clock
stops mattering; what survives is the universal $2\varepsilon$ behaviour, and the
two models agree to leading order. They do not agree at the next order: the
continuous constant has the elementary expansion
$\Ac = 2\varepsilon(1-\varepsilon+O(\varepsilon^2))$, whereas
\cref{a:thm:nearcritical} shows the discrete correction carries a factor
$\log(1/\varepsilon)$ and is therefore asymptotically larger.

\subsection*{Where the closed form went}

In continuous time the survival probability is elementary at every
time $t$, not merely asymptotically:
\begin{equation}
\label{a:eq:ctsSagain}
S(t) = \frac{\mu-\lambda}{\mu-\lambda\ee^{-(\mu-\lambda)t}}\cdot\frac{\mu}{\mu}
\quad\text{rearranges to}\quad
S(t) = \frac{\mu-\lambda}{\mu}\,\ee^{-(\mu-\lambda)t}
\lrb{1-\frac{\lambda}{\mu}\ee^{-(\mu-\lambda)t}}^{-1},
\end{equation}
so the amplitude $(\mu-\lambda)/\mu$ can simply be read off the exact solution.
In discrete time no such exact solution exists: the survival probability is defined
by iterating a quadratic map, and its amplitude is the infinite product
\cref{a:eq:prodFormula} of every early correction. The continuous process is
solvable because its generating function satisfies a linear first-order partial
differential equation with an elementary characteristic; the discrete one is not
because iterating a quadratic map is, in the precise sense of \cref{a:sec:ht}, a
transcendental operation.
